\documentclass[titlepage,a4paper]{jsarticle}
\usepackage{../../../sty/import}% 各種パッケージインポート
\usepackage{../../../sty/title}% タイトルページの変更
\usepackage{../../../sty/answergrid}
%% タイトルページの変数
% レポートタイトル
\title{心理学特論第6回目出席票}
% 提出日
\expdate{\today}
% 科目名
\subject{心理学特論}
% 分野
\class{情報経営システム工学分野}
% 学年
\grade{M1}
% 学籍番号
\mynumber{24336488}
% 記述者
\author{本間三暉}

\begin{document}
% titleページ作成
\maketitle
\section*{keyword\#1}
端数効果
\section*{keyword\#2}
心理的リアクタンス
\section{心理学特論第06回で取り上げたブランドとマーケティング，心理的テクニックについて，今後の就職活動や研究活動などでも活用できるように簡潔にまとめてみましょう！}

\subsection{自分の好きなブランドがあれば，なぜそれに魅力を感じるのか，簡潔に説明して下さい．}
私の好きなブランドはAppleである．魅力を感じる理由は，製品のデザインがシンプルで洗練されており，直感的に操作できる点にある．また，iPhoneやMac，iPadなどの製品同士がスムーズに連携するため，作業効率が高まる．
さらに，新しい技術を積極的に取り入れ，利用者にとって使いやすい製品を提供し続けている点にも魅力を感じる．
\subsection{今回取り上げたいくつかの心理的テクニックで，今後意識的に実践してみたいと思ったもの，
過去に自分が実践した（または他人から実践された）経験があるものについて，理由や状況を挙げて簡潔に書いて下さい．思いつく経験がなければ，単純な感想でも結構です．}
今回の授業で紹介された心理的テクニックの中で，今後意識的に実践してみたいと思ったものは「メタ認知」である．
メタ認知とは，自分の思考や感情を一段高い視点から客観的に観察し，必要に応じて行動を修正することである．

私はゲームや勉強でうまくいかないときに，「自分には才能がない」と考えてしまうことがある．
しかし，そのようなときに自分の考え方を客観的に見直すことで，「努力の方法に改善の余地があるのではないか」と考えられるようになると思う．
また，感情的になっているときにも，一度立ち止まって自分の状態を確認することで，冷静に行動できると感じた．

認知バイアスは誰にでも生じるものであり，それを完全になくすことは難しいが，メタ認知を意識することで，より適切な判断ができるようになると感じた．
\newpage
\section{心理学特論課題サイコロワーク解答記入欄}
\AnswerGridFilled{
  {認知的ケチ}
  {同調}
  {キティ・ジェノヴィーズ事件}
  {傍観者効果}
  {30％}
  {互恵性規範}
  {チャルディーニ}
  {譲歩的要請法}
  {フリードマン}
  {段階的要請法}
  {セルフ・ハンディキャッピング}
  {バーグラス}
  {栄光浴}
  {自己呈示}
  {身体像境界}
  {イーブン・ア・ペニー}
  {ザッツ・ノット・オール}
  {ロー・ボール}
  {ランチョン}
  {インフォームドコンセント}
}

\end{document}